%!TeX root=csp-audit.tex
\section{The register}\label{sec:register}

Nineteen places where \cite{ZhukSimplified} is wrong, incomplete, or states
less than its own proof gives. Sixteen are repaired and none of those repairs
needs new mathematics. Three are open: $\mathrm{D}17$, $\mathrm{D}19$ and
$\mathrm{D}21$ --- though none of the first two now blocks the proof paper in
its original form, since the route-four rewrite of
Entry~\ref{ent:route-four} moved what remains of both into one claim and one
stated lemma there; and $\mathrm{D}21$'s live repair route folds it into the
same claim (\S\ref{sec:D21}).

\begin{center}\footnotesize
\begin{longtable}{@{}llp{0.38\textwidth}p{0.26\textwidth}@{}}
\hline
\textbf{Tag} & \textbf{Where} & \textbf{What} & \textbf{Proof paper} \\
\hline\endhead
$\mathrm D1$ & \S2.3 & Lemma~19 restates a corollary of \cite{ZhukStrong}
  dropping two hypotheses; false as printed & \pkey{lem:ba-center-implies} \\
$\mathrm D2$ & \texttt{main:1682} & Propagation modulo a congruence: stated,
  used five times, never proved & \pkey{lem:propagation} \\
$\mathrm D4$ & \texttt{SS:2594} & $n=2$ for type $\TC$ in stable intersection
  is argued only for $\TPC$ & \pkey{thm:stable-intersection} \\
$\mathrm D6$ & \S3.2 & Bridges-from-relations applied without its third
  hypothesis, which can fail & \pkey{lem:connected} \\
$\mathrm D7$ & \S3.3 & Two gaps in the main induction: \emph{connected} proved
  where \emph{linked} is demanded, and (1c) derived for the wrong reduction &
  \pkey{thm:main-inductive} \\
$\mathrm D8$ & \texttt{main:3526} & Minimal-containing-a-point is used where
  minimal-outright is required; the bridging lemma was deleted &
  \pkey{lem:minimal-consistent} \\
$\mathrm D{11}$ & \texttt{SS:1154} & The symmetrisation formula defines a
  relation that is never a bridge & \pkey{lem:symmetrise} \\
$\mathrm D{12}$ & \texttt{SS:1103} & Bridges and abelian groups: \emph{false as
  stated}, with a three-element counterexample &
  \pkey{lem:nice-bridge-abelian} \\
$\mathrm D{13}$ & \texttt{SS:1273} & Linkedness asserted from an inclusion that
  bears on a different hypothesis & \pkey{lem:nontrivial-bridge} \\
$\mathrm D{14}$ & \texttt{SS:1394} & The box form of
  absorption-preserves-linkedness cited for a restriction that is not a box &
  \pkey{lem:absorb-linked-sub} \\
$\mathrm D{15}$ & \texttt{SS:2208} & Existence of an irreducible congruence
  factors by the wrong congruence, hiding an unsupported properness step &
  \pkey{lem:main-existence} \\
$\mathrm D{16}$ & \texttt{SS:2743} & Linear congruences on top: the conclusion
  is asserted & \pkey{lem:linear-on-top} \\
$\mathrm D{17}$ & \texttt{SS:3113} & Maximal multi-type extension concludes
  that two sets meet because each meets a third. \textbf{Open}, but the lemma
  is no longer consumed (Entry~\ref{ent:route-four}) &
  \pkey{lem:maximal-mult} \\
$\mathrm C{10}$ & \texttt{SS:208} & The central-relation citation drops the
  third alternative of the theorem it cites &
  \pkey{lem:central-relation} \\
$\mathrm F$ & \texttt{SS:545, 837, 921} & A full fibre inferred from a
  surjective projection, at three sites & \pkey{lem:full-fibre} \\
$\mathrm D{18}$ & \texttt{main:1373} & The definition of a linear congruence
  is ill-typed: it intersects a relation on elements with a set of tuples of
  classes & \pkey{def:linear-congruence} \\
$\mathrm D{19}$ & \texttt{main:2663, 2693} & Two stable-intersection
  applications in the typing argument never establish the corollary's
  leave-one-out hypothesis, and it is not recoverable at coordinates where
  $D^{(2)}$ equals $D^{(1)}$. \textbf{Partly repaired}: the anchored form
  now exists and settles the mixed-type corner; the same-class corner is
  open &
  \pkey{lem:anchored-\brk intersection}, \pkey{cl:typing-anchor} \\
$\mathrm D{20}$ & \texttt{main:2574} & The parallelogram lemma omits the
  subdirectness its own proof spends, and the repair needs $\TS$-freeness of
  the $D^{(1)}$-domains besides &
  \pkey{lem:parallelogram-\brk crucial} \\
$\mathrm D{21}$ & \texttt{SS:2892} & The linkedness-preservation auxiliary
  asserts its cross-component step: the $S$-witness and the mixed edge may
  lie in different linked components, and nothing joins them.
  \textbf{Open}: the anchored form is proved, the unanchored lemma stays
  stated &
  \pkey{lem:preserve-linked-\brk anchored}, \pkey{lem:preserve-linked} \\
\hline
\end{longtable}
\end{center}

Two patterns recur rather than occurring once, and are worth naming because a
reader checking the proof paper will meet them repeatedly. The first is a
\emph{proper} containment asserted where the cited lemma gives only
$\subseteq$, usually after factoring by a congruence, which can close the gap;
it appears at least five times and is repaired the same way each time, by
producing the witness the source has already displayed two lines earlier. The
second is an upgrade from ``the image has size one'' to ``the image is
\emph{this} block''; it appears four times and the proof paper makes it a named
step.

\subsection{$\mathrm D1$ --- two hypotheses dropped, and the statement becomes false}
\label{sec:D1}

\cite[Lem.~19]{ZhukSimplified} restates \cite[Cor.~6.1.2, 6.9.2]{ZhukStrong}.
The source it cites reads ``Suppose $R\le A_{1}\times\dots\times A_{n}$,
$\proj_{1}(R)=A_{1}$, $B_{i}\le_{\TBA(t)}A_{i}$ for every $i$''. The
restatement drops the subdirectness $\proj_{1}(R)=A_{1}$, and drops --- in the
$\TBA$ case --- that the $B_{i}$ absorb with a \emph{common} term $t$, which is
the reason \cite{ZhukStrong} carries $t$ in the notation.

\emph{Without the first the statement is false.} Take
$R=\{(0,0)\}\le\mathbf Z_{p}\times\mathbf Z_{p}$, a subuniverse by
idempotence, and $C_{i}=A_{i}$, which is allowed because $\subte{T}$ admits
equality. Then $R\cap(C_{1}\times C_{2})=R$ and $\proj_{1}(R)=\{0\}$, which is
neither empty nor all of $\mathbf Z_{p}$; the conclusion therefore asserts
$\{0\}\subt{\TBA}\mathbf Z_{p}$, contradicting the paper's own Lemma~29.

\emph{Repair.} Restore both hypotheses. The proof paper states them in
\pkey{lem:ba-center-implies}. The common-term condition is not free at the
call site: \pkey{thm:main-inductive}(2) records only that every coordinate
reduction has type $\TBA$, and applies the lemma to the relation computing
$E_{2}$ from $E_{1}$. Either a lemma synthesising one common binary term from a
finite family of coordinate absorptions is needed, or the reduction data must
carry a common witness from the start. This is recorded as an obligation on the
proof paper, not as a defect of the source, since the source's own version of
the lemma carries the hypothesis.

\subsection{$\mathrm D2$ --- propagation modulo a congruence is never proved}
\label{sec:D2}

Stated at \texttt{main:1682}; used at \texttt{SS:2506, 2514, 2588, 3110} and
\texttt{main:2999}. Never proved, and not attributed to any external source.

The paper's own convention (\texttt{main:1058}) is \emph{``we always duplicate
statements if the proof appears in a later section''}, implemented by
\texttt{\textbackslash newtheorem*} duplicates. Checking all thirteen
statements of \S2.3 against that convention leaves exactly two not restated in
\S5: this corollary, and \texttt{LEMBACenterImplies}, which is declared an
external import and so is absent by design.

\emph{Repair.} It is Propagation specialised to the canonical surjection
$\pi\colon\mathbf A\twoheadrightarrow\mathbf A/\delta$, whose clauses (f),
(ft), (fs), (fm) are its four clauses. The identification is literal rather
than approximate: $X/\delta$ is \emph{defined} as the image
$\{x/\delta:x\in X\}$, which is $\pi(X)$; in particular the corollary's
hypothesis ``$B/\delta$ is $\TS$-free'' and Propagation's ``$f(B)$ is
$\TS$-free'' are the same condition. Proof paper: \pkey{cor:prop-quotient},
derived from \pkey{lem:propagation}.

A mechanical audit of all $81$ labelled statements of the source
(\S\ref{sec:machine-checks}) shows this is the \emph{only} statement that is
cited while being neither proved nor attributed. That matters beyond this
entry: the audit covers \S5.6 and \S5.7, so whatever else is wrong there, it is
not another statement asserted from nowhere.

\subsection{$\mathrm D4$ --- $n=2$ for type $\TC$ is unargued}\label{sec:D4}

The stable-intersection theorem concludes, in its case (c), that $n=2$ and
$T_{1}=T_{2}=\TC$. The proof at \texttt{SS:2594}\,ff.\ reduces general $n$ to a
pair, obtains $T_{1}=T_{2}$, and then:

\begin{quote}\small
``If $T_{1}=\TPC$ then $\sigma_{1}\circ\sigma_{2}=\sigma_{1}=\sigma_{2}$ \dots
Additionally, this implies that $n$ cannot be greater than 2 for
$T_{1}=\TPC$, as in this case the intersection $C_{1}\cap C_{2}$ must be
empty.''
\end{quote}

That argument is specific to $\TPC$: equal congruences make $C_{1}$ and $C_{2}$
distinct blocks, hence \emph{pairwise} disjoint, which contradicts hypothesis
(3) once $n\ge3$. For type $\TC$ there is no congruence and no analogue, and
the proof offers nothing; hypothesis (2) gives only
$C_{1}\cap\dots\cap C_{n}=\varnothing$, never $C_{1}\cap C_{2}=\varnothing$.
The same gap is visible at \texttt{SS:2650}, where the $\TPC$ argument needs
one further unstated step --- the two $\sigma$-blocks are distinct because
$B_{1}\cap C_{2}'\neq\varnothing$ would otherwise put a point in
$C_{1}\cap C_{2}'$ --- which again has no type-$\TC$ analogue.

\emph{Repair.} The missing argument is \cite[Thm.~3.7]{ZhukStrong}, which is
the same statement for subuniverses of a \emph{single} algebra, resting on
\cite[Lem.~6.11]{ZhukStrong}: for $n\ge3$ and $C_{i}$ central in $A_{i}$, no
$(C_{1},\dots,C_{n})$-essential relation $R\le A_{1}\times\dots\times A_{n}$
exists. The 2024 setting is not literally that of Theorem~3.7 --- there the
$B_{i}$ are subuniverses of one algebra, whereas here $C_{i}$ is central in
$B_{i}$ and the $B_{i}$ differ --- so a reduction is needed first: put
$D=\bigcap_{i}B_{i}$ and $E_{i}=C_{i}\cap D$; then $D$ is a nonempty
subuniverse, each $E_{i}$ is properly central in $D$, $\bigcap_{i}E_{i}
=\varnothing$ and $\bigcap_{i\neq j}E_{i}\neq\varnothing$ for every $j$. The
diagonal of $D^{n}$ is then $(E_{1},\dots,E_{n})$-essential. Proof paper:
inside \pkey{thm:stable-intersection}.

\emph{Cost of the repair.} \cite[Lem.~6.11]{ZhukStrong} has a printed proof
that cites \emph{itself} three times where it means the preceding lemma, so a
self-contained rendering must supply both. Both citations were re-read verbatim
(Lemma~6.11 at 2005.00593 p.~37, Theorem~3.7 at p.~46).

\subsection{$\mathrm D6$ --- a hypothesis that can fail, and cruciality supplies it}
\label{sec:D6}

\texttt{LEMBridgeFromRelation} requires, besides subdirectness and
rectangularity of the first and last variables, that there exist
$(b_{1},\bar a,a_{n})$ and $(a_{1},\bar a,b_{n})$ in $R$ with
$(a_{1},\bar a,a_{n})\notin R$. The proof of \texttt{LEMConnectedProperties}(a)
invokes it for every constraint along a path and never discharges that
hypothesis.

\emph{It can fail.} Machine-checked witness:
$R=\{(x,z,y)\in\mathbb Z_{4}\times\mathbb Z_{2}\times\mathbb Z_{4}:
x\equiv z\equiv y\ (\mathrm{mod}\ 2)\}$ is subdirect, $\ConOne(R,1)$ is
congruence mod $2$ and coordinate $1$ is rectangular, and the third hypothesis
is unsatisfiable. The construction then yields
$\delta(x_{1},x_{2},y_{1},y_{2})=\exists z\,R(x_{1},z,y_{1})\wedge
R(x_{2},z,y_{2})$, whose projection onto the first two coordinates
\emph{equals} $\ConOne(R,1)$ rather than properly containing it --- so $\delta$
is not a bridge, clause (c) failing. The lemma is genuinely inapplicable and
genuinely necessary.

\emph{Repair.} Cruciality supplies the hypothesis. If the third hypothesis
fails at coordinates $1$ and $n$ then $R$ is the conjunction of its two
projections,
\[
  R=\{(x,\bar a,y):(x,\bar a)\in\proj_{[n-1]}(R)\ \text{and}\
     (\bar a,y)\in\proj_{[n]\setminus\{1\}}(R)\},
\]
the inclusion $\supseteq$ being exactly the failure of the hypothesis. A
constraint that is crucial in a reduction therefore satisfies the hypothesis at
every pair of its coordinates: otherwise weakening it replaces it by
constraints whose conjunction defines the same relation, so the weakened
instance has the same solution set and in particular still has no solution in
the reduction, contradicting cruciality. On the witness above,
$\proj_{1,2}(R)$ and $\proj_{2,3}(R)$ do conjoin to $R$ --- exactly the
degenerate shape the corollary rules out.

Proof paper: \pkey{lem:connected} is stated for crucial instances. This is a
notion \cite{ZhukSimplified} inherits from \texttt{LEMCrucialMeansIrreducible}
and then stops using.

\subsection{$\mathrm D7$ --- two gaps in the main induction}\label{sec:D7}

\paragraph{(i) The endgame proves \emph{connected}; (1c) demands \emph{linked}.}
Both endgames of Case~2 finish with ``\dots which means that $\mathcal I$ is
connected and satisfies (1b) if its solution set is subdirect or (1c)
otherwise''. But (1c) requires a \emph{linked} connected subinstance, and only
connectedness was established.

\emph{Repair, via irreducibility.} Suppose $\mathcal I$ is connected with
non-subdirect solution set. Take $\mathcal J=\Upsilon=\mathcal I$, legitimate
since an instance is a weakening, hence an expanded covering, of itself. It
remains to see $\mathcal I$ is linked. If it were not, $\mathcal I$ would
witness the failure of its own irreducibility: its variables are among its own,
each constraint is the projection of itself onto all its variables, it is not
fragmented (connectedness of the constraint-adjacency graph forces constraints
to share variables), it is not linked, and its solution set is not subdirect
--- precisely conditions (i)--(v) of the definition. Since $\mathcal I$ is
irreducible by standing hypothesis, it is linked. Two lines, and they explain
why irreducibility is a hypothesis of the theorem at all.

\paragraph{(ii) Case~1 derives (1c) for the wrong reduction.}
Case~1 obtains a $1$-consistent $D^{(2)}\subte{T}D^{(1)}$, shows $\mathcal I$
is crucial in $D^{(2)}$, and closes with ``applying the inductive assumption to
$D^{(2)}$ we derive the required conditions''. Conditions (1a) and (1b) mention
no reduction and transfer verbatim. (1c) does not: it asserts an expanded
covering \emph{crucial in the ambient reduction}. The inductive hypothesis
delivers $\mathcal K$ crucial in $D^{(2)}$; the goal needs a covering crucial
in $D^{(1)}$, and the containment runs the wrong way --- an instance with no
solution in the smaller $D^{(2)}$ may have one in $D^{(1)}$.

\emph{First, the induction structure, which is load-bearing and unstated.} The
proof opens ``We prove the claim by induction on the size of $D^{(1)}$.
\emph{Let us prove (2) first.}'' That ordering is not stylistic. Part~(2) at a
given measure uses (1) only at strictly smaller measure, whereas part~(1) at
that measure uses \textbf{(2) at the same measure} --- Case~1 does exactly this
at \texttt{main:3507}. So the induction is: strong induction on
$k=\mu(D^{(1)})$; at each level establish (2) for \emph{all} instances, then
(1) for all instances. That is acyclic: $(2)_{k}\leftarrow(1)_{<k}$,
$(1)_{k}\leftarrow(2)_{k},(1)_{<k}$.

\emph{Why that alone does not close the gap.} One would like to apply (2) to
$\mathcal K$ and conclude contrapositively that $\mathcal K^{(1)}$ has no
solution. But $\mathcal K$ is an expanded \emph{covering}, so it has more
variables than $\mathcal I$, and under $\mu=\sum_{x}|D^{(1)}_{x}|$ its measure
is \emph{larger}, not equal. The device that works for a weakening does not
work for a covering.

\emph{The repair, which avoids re-entering the induction.} Put
$\mathcal L=\mathcal K\wedge\mathcal I$. It is an expanded covering of
$\mathcal I$, since $\mathcal K$ is one, $\mathcal I$ is a weakening of itself
hence one, and the union of two expanded coverings is one. $\mathcal L^{(1)}$
has no solution, since $\mathcal L$ contains every constraint of $\mathcal I$
and $\mathcal I$ is crucial in $D^{(1)}$. So $\mathcal L$ may be weakened to an
instance $\mathcal L'$ crucial in $D^{(1)}$. Every constraint of $\mathcal K$
survives that weakening: each is crucial in $D^{(2)}$, so weakening it yields a
solution in $D^{(2)}\subseteq D^{(1)}$, and weakening it inside any weaker
instance yields only more solutions; so each is already crucial in $D^{(1)}$ at
every stage, and the process only weakens constraints that are not. Therefore
$\Upsilon\subseteq\mathcal K\subseteq\mathcal L'$, its properties being
intrinsic to $\Upsilon$ and untouched, and $\mathcal L'$ is an expanded
covering of $\mathcal I$ crucial in $D^{(1)}$. No appeal to the induction at a
larger measure, so the measure survives intact.

Proof paper: \pkey{thm:main-inductive} and \pkey{haz:case1c}.

\begin{entry}[The measure, and a rejected objection]\label{ent:measure}
$\mu(D)=\sum_{x}|D_{x}|$ is well founded for every appeal the induction makes.
An earlier draft of this entry, and the external review independently, held
that it is not --- on the ground that the induction is applied at expanded
coverings, where the variable set grows and $\mu$ grows with it. \textbf{That
is wrong, and the enumeration settles it.}

The proof of \texttt{THMMainInductiveCSPClaim} runs from \texttt{main:3431} to
\texttt{main:3858} and makes exactly seven appeals to its own inductive
assumption: \texttt{main:3439}, \texttt{3507}, \texttt{3511}, \texttt{3536},
\texttt{3689}, \texttt{3714}, \texttt{3724}. Each is to a \emph{weakening} of
the current instance --- so $\Var$ can only shrink --- with a reduction
strictly smaller in $\mu$: $D^{(2)}$ in the first four, and $D^{(B,\bot)}$ in
the last three, the latter being ``the maximal $1$-consistent (probably empty)
reduction \emph{for $\mathcal I$}'' below $D^{(B,\top)}$
(\texttt{main:3554}), which cuts $D^{(1)}_{z}$ to a proper
$B<^{D_{z}}_{T(\sigma)}D^{(1)}_{z}$. The three in the
$\mathcal B_{0}\neq\varnothing$ branch draw their instances from $\Omega$,
whose \emph{condition 1} (\texttt{main:3629}) is ``every instance in $\Omega$
is a weakening of $\mathcal I$''. Expanded coverings occur in that proof only
as outputs --- alternative (1c), and the auxiliary $\bot(\cdot)$,
$\Delta(\cdot)$ and $\Theta$ --- which carry no well-foundedness obligation.

What misled both readings is \texttt{main:3913}, ``By Theorem
\texttt{THMMainInductiveCSPClaim} applied to $\mathcal J'$ and
$D^{(G,\bot)}$'', where $\mathcal J'$ is a weakening of an expanded covering.
That line is at \texttt{main:3913}, inside \texttt{THMPCDoesnotKillAllSolutions}
(\texttt{main:3860--3976}), a \emph{separate} theorem proved afterwards and
consuming the main claim as a black box; the main proof does not cite it back,
so there is no mutual induction. Applying a proved theorem to a covering is
free. Proof paper: \pkey{haz:the-induction}(c).
\end{entry}

\subsection{$\mathrm D8$ --- minimal-containing-a-point versus minimal-outright}
\label{sec:D8}

At \texttt{main:3526} the main induction picks, for each variable, ``the
minimal $D^{(2)}_{x}\le_{\MT T}D^{(1)}$ containing $s(x)$'', and two lines later
invokes a lemma whose hypothesis (5) requires $D^{(2)}_{x}$ to be \emph{a
minimal $\MT T$ subuniverse} --- minimal outright. Those are not the same
condition. The lemma bridging them has been deleted, along with its citation:

\begin{quote}\footnotesize\ttfamily
3526\quad For every variable x choose the minimal D\^{}(2)\_x <=\_\{MT\} D\^{}(1) containing s(x).\\
3527\quad \%By Lemma \textbackslash ref\{LEMMinimalContainingIsMinimal\}\\
3528\quad \%\$D\^{}\{(2)\}\_\{x\}\$ is a minimal \textbackslash mathcal\{M\}T subuniverse.\\
3529\quad By Lemma \textbackslash ref\{LEMMultiTypeStillStable\}
\end{quote}

The deleted statement (\texttt{main:1864}) was: if $C$ is inclusion-minimal with
$b\in C$ and $C\le_{\MT T}B$, then no $C'\subsetneq C$ has $C'\le_{\MT T}B$.

\emph{Repair.} Fix $T\in\{\TL,\TPC,\TD\}$ and $B\lll A$, let $\Sigma$ be the
finite set of congruences that are $T$-dividing for $B$, and for $b\in B$ put
$M(b)=B\cap\bigcap_{\sigma\in\Sigma}[b]_{\sigma}$, discarding any $\sigma$ with
$B\cap[b]_{\sigma}=B$. Then $M(b)$ is the least $\MT T$ subuniverse of $B$
containing $b$; and if $c\in M(b)$ then $[c]_{\sigma}=[b]_{\sigma}$ for every
$\sigma\in\Sigma$, so $M(c)=M(b)$. Consequently an inclusion-minimal $\MT T$
set containing $b$ equals $M(b)$ and admits no proper $\MT T$ subset: any such
$C'$ is nonempty, and any $c\in C'$ has $M(c)=M(b)\subseteq C'$.

The content is the second sentence: $b\mapsto M(b)$ is a closure whose image
consists of pairwise disjoint minimal sets, because membership in $M(b)$ is
decided blockwise and blocks of an equivalence relation either coincide or are
disjoint. Proof paper: \pkey{lem:minimal-consistent} and
\pkey{haz:minimality}.

\subsection{$\mathrm D{11}$--$\mathrm D{13}$ --- three defects, one missing hypothesis}
\label{sec:D11}

\S5.4 of \cite{ZhukSimplified} is the only part of \S5 that manipulates bridges
as objects in their own right rather than consuming them, and all three of its
defects are about which auxiliary property survives a construction. All three
are cured by one hypothesis, which the source already has in its vocabulary:
\texttt{LEMPCCongruencePropertyInductiveStep} (\texttt{SS:1379}) carries, as
hypothesis (2), exactly \emph{pair-reflexivity} --- $(a,b,a,b),(b,a,b,a)\in
\delta$ for every $(a,b)\in\proj_{1,2}(\delta)$.

\paragraph{$\mathrm D{11}$: the symmetrisation formula.}
The proof of \texttt{LEMBlockOf\brk GoodBridge\brk DoesNotHaveBAC} opens by symmetrising $\delta$:
\begin{quote}\footnotesize\ttfamily
1154\quad \textbackslash sigma(x\_1,x\_2,x\_3,x\_4) = \textbackslash exists x\_5 \textbackslash exists x\_6\\
1155\quad \textbackslash delta(x\_1,x\_2,x\_5,x\_5)\\
1157\quad \textbackslash wedge \textbackslash delta(x\_3,x\_4,x\_5,x\_6).
\end{quote}
Clause (d) of the bridge definition (\texttt{main:1258}) is
``$(a_{1},a_{2},a_{3},a_{4})\in\delta$ implies
$(a_{1},a_{2})\in\sigma_{1}\iff(a_{3},a_{4})\in\sigma_{2}$''. The first
conjunct has $x_{5}$ in both of the last two places, so $(x_{5},x_{5})\in 0_{A}$
forces $x_{1}=x_{2}$. Hence $\proj_{1,2}(\sigma)\subseteq 0_{A}$, while clause
(c) demands $\proj_{1,2}(\sigma)\supsetneq\sigma_{1}=0_{A}$: \emph{the relation
defined is never a bridge}, for any $\sigma_{1}$, and every later line of the
proof is false of it.

It is a one-character slip. The intended relation is $\delta\ast\delta^{-1}$,
and the source writes exactly that, correctly, $210$ lines later at
\texttt{SS:1365--1367}, for the same purpose.

\paragraph{$\mathrm D{12}$: bridges and abelian groups, false as stated.}
\texttt{LEMNiceBridgeGivesAbelianGroup} (\texttt{SS:1103}): for $\sigma$ a
congruence and $\delta$ a bridge from $\sigma$ to $\sigma$ with
$\proj_{1,2}(\delta)=\bcol\delta=A^{2}$ and $\delta$ symmetric, there is an
abelian group $G$ with $(A/\sigma;\delta/\sigma)\cong(G;x_{1}-x_{2}=x_{3}-x_{4})$.

\emph{Counterexample.} $\mathbf A=\mathbf Z_{3}$ with
$w(x_{1},\dots,x_{4})=x_{1}+x_{2}+x_{3}+x_{4}\bmod 3$ --- idempotent since
$4\equiv 1$, and a special WNU, so $\mathbf A\in\mathcal V_{4}$. Take
$\sigma=0_{\mathbf A}$ and
$\delta=\{(x_{1},x_{2},x_{3},x_{4})\in\mathbb Z_{3}^{4}:
x_{1}-x_{2}+x_{3}-x_{4}=0\}$. Verified by exhaustion: $\delta$ is a subuniverse
of $A^{4}$; clause (d) holds; $\proj_{1,2}(\delta)=\bcol\delta=A^{2}$; $\delta$
is symmetric; and none of the six bijections $\mathbb Z_{3}\to\mathbb Z_{3}$
carries $\delta$ to $\{x_{1}-x_{2}=x_{3}-x_{4}\}$. Since every abelian group of
order $3$ is $\mathbb Z_{3}$, no $G$ works.

\emph{What is true.} $\delta$ contains all $(a,a,b,b)$, so it is the preimage
of some $K\le G\times G$ under $(x_{1},x_{2},x_{3},x_{4})\mapsto
(x_{1}-x_{2},x_{3}-x_{4})$. Clause (d) makes $K$ meet $G\times 0$ and
$0\times G$ only at the origin, and subdirectness makes it a graph, so
$\delta=\{\varphi(x_{1}-x_{2})=x_{3}-x_{4}\}$ for an automorphism $\varphi$.
Symmetry says $\varphi^{2}=\mathrm{id}$; Zhuk's conclusion is the case
$\varphi=\mathrm{id}$, and the counterexample is $\varphi=-\mathrm{id}$, which
is why it needs $p$ odd. The step that breaks is ``composing the bridge
$\delta$ with itself we get a bridge $\delta_{0}\supseteq\delta$'':
$\delta\ast\delta=\{\varphi^{2}(x_{1}-x_{2})=x_{3}-x_{4}\}$, which for
$\varphi=-\mathrm{id}$ does not contain $\delta$, and the closing count
$|\delta|=|\delta_{0}|=|A|^{3}$ needs the containment it does not have.

\paragraph{$\mathrm D{13}$: linkedness is asserted, not shown.}
\texttt{SS:1273} puts $\delta'=\delta\cap B^{4}\cap(\cover\sigma\times
\cover\sigma)$ for $B$ a block of $\operatorname{LeftLinked}(\cover\sigma)$ and
says ``Since $\bcol\delta\supseteq\cover\sigma$, $\delta'$ satisfies all the
conditions of Lemma \texttt{LEMBlockOfGoodBridgeDoesNotHaveBAC}''. Two of those
conditions are that $\proj_{1,2}(\delta')$ properly contains $0_{B}$ and is
\emph{linked}. What $\bcol\delta\supseteq\cover\sigma$ gives is the third
condition. Nothing offered bears on linkedness --- and it does not follow: the
witness $(x_{3},x_{4})$ supplied by minimality of $\cover\sigma$ lies in
$\cover\sigma$ but in no particular block, and if it lands outside $B$ for
every choice then $\proj_{1,2}(\delta')$ collapses to $0_{B}$ and $\delta'$ is
not even a bridge. The same gap recurs at \texttt{SS:1290}.

\paragraph{The single cure, and why it costs nothing.}
Pair-reflexivity kills all three. $\mathrm D{13}$: $\cover\sigma\subseteq
\proj_{1,2}(\delta)$ by minimality of $\cover\sigma$, since $\proj_{1,2}(\delta)$
is a subalgebra of $A^{2}$, stable under $\sigma$ by clause (a) and properly
containing $\sigma$ by clause (c); so pair-reflexivity puts
$(x_{1},x_{2},x_{1},x_{2})$ in $\delta'$ for every
$(x_{1},x_{2})\in\cover\sigma\cap B^{2}$, giving $\proj_{1,2}(\delta')=
\cover\sigma\cap B^{2}$ exactly --- which is linked precisely because $B$ is a
block. $\mathrm D{12}$: pair-reflexivity \emph{is} $\varphi=\mathrm{id}$, and it
also repairs the proof directly, since $\delta\subseteq\delta\ast\delta$ by
taking $(x_{1},x_{2})$ itself as the intermediate pair. $\mathrm D{11}$:
$\delta\ast\delta^{-1}$ is pair-reflexive by construction, so the corrected
formula delivers what the lemma needs.

And it costs nothing, because the two lemmas that need the hypothesis are each
invoked exactly once in the whole paper --- \texttt{LEMNontrivialReflexive\-Bridge\-Implies}
at \texttt{SS:1368}, on a $\delta\ast\delta^{-1}$, and
\texttt{LEMNiceBridgeGivesAbelianGroup} from inside that call.

Proof paper: \pkey{def:pair-reflexive}, \pkey{lem:symmetrise},
\pkey{lem:nice-bridge-abelian}, \pkey{lem:nontrivial-bridge},
\pkey{rem:phi-twist}.

\subsection{$\mathrm D{14}$ --- the box form cannot reach the subrelation}
\label{sec:D14}

\texttt{SS:1394}, inside Case~1 of
\texttt{LEMPCCongruencePropertyInductiveStep}: with $B\subt{T}A$, $|B|>1$ and
$\delta'=\delta\cap B^{4}$,

\begin{quote}\small
``By Lemma \texttt{LEMBACenterLinkedness}
$\operatorname{RightLinked}(\proj_{1,2,3}(\delta'))=B^{2}$.''
\end{quote}

\texttt{LEMBACenterLinkedness} (\cite[Prop.~2.15(i)]{BartoKozik}) restricts a
binary relation by a \emph{box} $B_{1}\times B_{2}$. Applied with
$W=\proj_{1,2,3}(\delta)$, $B_{1}=P\cap B^{2}$, $B_{2}=B$ it yields linkedness
of $\proj_{1,2,3}(\delta)\cap B^{3}$, which contains
$\proj_{1,2,3}(\delta\cap B^{4})$ and can be strictly larger: a tuple of
$\delta$ whose first three coordinates lie in $B$ need not have its fourth
there. The two really do differ --- for the meet-semilattice on $\{0,1,2\}$
with $B=\{0,1\}$, the subalgebra $K=\{(0,0),(1,2)\}$ of $A^{2}$ has
$\proj_{1}(K)\cap B=\{0,1\}$ but $\proj_{1}(K\cap B^{2})=\{0\}$. Getting
$\delta\cap B^{4}$ out of the box form instead would need linkedness of
$\delta$ in the split $(1,2)\mid(3,4)$ or $(1,2,3)\mid(4)$, and hypothesis (3)
gives neither: a path in the coarser bipartite graph does not lift, because
consecutive edges through a third coordinate may use different fourth
coordinates. The lemma's own hypothesis $R\cap(B_{1}\times B_{2})\sdle
B_{1}\times B_{2}$ is also never discharged.

\emph{The assertion is true; the citation is to the wrong form of the
principle.} What closes it is the \emph{subrelation} form: if
$W\sdle\mathbf P\times\mathbf A$ is linked and $\varnothing\neq W'\subte{T}W$
with $T\in\{\TBA,\TC\}$, then $W'$ is linked over its own projections. That
follows from tools already present in \S5 --- pp-propagation, intersection, and
factoring --- together with one elementary fact the paper never states:
$0_{\mathbf C}$ is not an absorbing subuniverse of $\mathbf C^{2}$ when
$|C|\ge2$, because absorption at position $i$ forces the term not to depend on
argument $i$, at every $i$, hence to be constant, and idempotence then forces
$|C|=1$.

The subrelation form also discharges, for free, the subdirectness hypothesis
the box form needed. Proof paper: \pkey{lem:no-absorbing-diagonal},
\pkey{lem:absorb-linked-sub}, \pkey{haz:box-vs-sub}.

\emph{Contrast.} The analogous step in \texttt{LEMPCBridgesAreTrivial} Case~2
\emph{is} justified by the source --- by first proving the closure property
$\xi\cap(A^{2}\times B\times A)\subseteq B^{4}$ (\texttt{SS:1548--1554}), which
makes $\proj_{1,2,3}(\xi\cap B^{4})=\proj_{1,2,3}(\xi)\cap(A^{2}\times B)$.
Here $B$ is a strong subuniverse rather than a linkedness block, no closure
property is available, and absorption has to do the work instead.

\subsection{$\mathrm D{15}$ --- read $\delta$ for $\sigma$, and restore stability}
\label{sec:D15}

\texttt{SS:2208--2224}. A congruence $\delta\supseteq\sigma$ is chosen maximal
with $|B/\delta|>1$, and $B/\delta$ is shown BA and centre free. To show
$\delta$ irreducible the proof supposes $\delta=S_{1}\cap\dots\cap S_{k}$ with
$S_{i}\supsetneq\delta$, finds an $i$ with $B^{2}\not\subseteq S_{i}$, and then:

\begin{quote}\small
``By Lemma \texttt{LEMLeftLinkedStayFull}
$\operatorname{LeftLinked}(S_{i}\cap(B\times B))/\sigma=(B/\sigma)^{2}$. Hence
$(S_{i}\cap(B\times B))/\sigma$ is a linked relation, which by Lemma
\texttt{LEMLinkedImpliesBACenter} implies the existence of a BA or central
subuniverse on $B/\delta$.''
\end{quote}

Two things are wrong, and they are the same thing twice. The relation is
factored by $\sigma$ while the conclusion is about $B/\delta$; those do not
match, since $(S_{i}\cap B^{2})/\sigma$ lives on $B/\sigma$. And properness is
not established: \texttt{LEMLinkedImpliesBACenter} needs a \emph{proper}
subdirect relation, $B^{2}\not\subseteq S_{i}$ gives only
$S_{i}\cap B^{2}\subsetneq B^{2}$, and factoring can close that gap.

\emph{Repair.} Read $\delta$ for $\sigma$ throughout the last four lines, which
is plainly what was meant --- $B/\delta$ was shown BA and centre free two lines
earlier for exactly this use, and the stated conclusion already says
$B/\delta$. Properness then follows, but only from a clause the proof drops
when it introduces the $S_{i}$: the definition of irreducible
(\texttt{main:1237}) requires the $S_{i}$ to be \textbf{stable under $\delta$}.
Stability makes $S_{i}$ a union of products of $\delta$-blocks, so
$(S_{i}\cap B^{2})/\delta=(B/\delta)^{2}$ would force $B^{2}\subseteq S_{i}$.
That clause is the only part of the definition of irreducibility doing work
here. Proof paper: \pkey{lem:main-existence}, \pkey{haz:main-existence}.

This propagates: \texttt{LEMUbiquity} rests on this lemma, so $\mathrm D{15}$
reaches one of the headline properties.

\subsection{$\mathrm D{16}$ --- the conclusion is the whole lemma, and is asserted}
\label{sec:D16}

\texttt{SS:2743}: $\sigma$ linear on $\mathbf A\in\Vn$ with
$\cover\sigma=A^{2}$ implies $\mathbf A/\sigma\cong\mathbf Z_{p}$. The proof
produces the perfect-linear witness $\zeta\le\mathbf A\times\mathbf A\times
\mathbf Z_{p}$ with $\proj_{1,2}(\zeta)=A^{2}$, sets $\xi(x,z)=\zeta(x,a,z)$,
and then says in full: ``Then $\xi$ is a bijective relation giving an
isomorphism $\mathbf A/\sigma\cong\mathbf Z_{p}$.'' That is the entire content
of the lemma. What $\zeta$ gives directly is only $\xi^{-1}(0)=[a]_{\sigma}$.

\emph{Repair, by a different route.} The lemma follows from a statement already
repaired here. A bridge $\delta$ from $\sigma$ to $\sigma$ with
$\bcol\delta\supsetneq\sigma$ exists by the linear criterion; replace it by
$\delta\ast\delta^{-1}$, symmetric and pair-reflexive
(\S\ref{sec:D11}). Then \texttt{LEMNontrivialReflexiveBridgeImplies}(3)
applies, and $\cover\sigma=A^{2}$ is a congruence with the single block $A$, so
$\mathbf A/\sigma\cong\mathbb Z_{p}^{m}$ as a group. Now $m=0$ gives
$\sigma=A^{2}$, impossible because an irreducible congruence is proper
(\S\ref{sec:leg-emptyfamily}); and $m\ge2$ contradicts irreducibility, because
$\mathbf A/\sigma$ is affine, so the coordinate congruences
$\theta_{i}=\{(x,y):x_{i}=y_{i}\}$ are linear subspaces of
$(\mathbb Z_{p}^{m})^{2}$, hence subuniverses of $(\mathbf A/\sigma)^{2}$, each
properly above $0_{\mathbf A/\sigma}$ with $\bigcap_{i}\theta_{i}=
0_{\mathbf A/\sigma}$. Hence $m=1$. That the group isomorphism is an
isomorphism \emph{of algebras} is the proposition of
\S\ref{sec:leg-Zp}, and without it the statement does not typecheck against
$\mathbf Z_{p}\in\Vn$.

This route uses only in-paper material, where the source's route imports
\cite[Cor.~8.17.1]{ZhukJACM} --- so the repair is \textbf{one fewer external
dependency} than the source. Proof paper: \pkey{lem:linear-on-top},
\pkey{haz:on-top}.

\begin{entry}[Salvage: the source's $\zeta$ really is a function]\label{ent:zeta}
Worth keeping, because it repairs half of the original argument and is needed
if the perfect-linear machinery is used elsewhere. Put
$t(x,y,z)=w(x,y,\dots,y,z)$ with $n-2$ copies of $y$; on $\mathbb Z_{p}$ this
is $x-y+z$, because $p\mid n-1$. Applying $t$ to three tuples of $\zeta$ over
the same $(a_{1},a_{2})$ gives, by idempotence, that the fibre is closed under
$x-y+z$, hence is a coset of a subgroup of $\mathbb Z_{p}$: empty, a singleton,
or all of $\mathbb Z_{p}$. It is not all of $\mathbb Z_{p}$, since that fibre
would contain $0$ and a nonzero element, making $(a_{1},a_{2})$ both
$\sigma$-related and not; and it is nonempty on $\cover\sigma$ since
$\proj_{1,2}(\zeta)=\cover\sigma$. So $\zeta$ is the graph of a homomorphism
$G\colon\mathbf A^{2}\to\mathbf Z_{p}$ with $G^{-1}(0)=\sigma$.
\end{entry}

\subsection{$\mathrm D{17}$ --- maximal multi-type extension. \textbf{Open}}
\label{sec:D17}

\texttt{SS:3113}. $\sigma$ is maximal with $(C_{1}\circ\sigma)\cap
B_{2}=\varnothing$; Case~1 supposes $E\subt{\TS}B_{1}/\sigma$ and argues:

\begin{quote}\small
``By Theorem \texttt{THMMainStableIntersection} $E\cap B_{2}/\sigma\neq
\varnothing$ and $E\cap C_{1}/\sigma\neq\varnothing$, hence
$C_{1}/\sigma\cap B_{2}/\sigma\neq\varnothing$, which contradicts our
assumptions.''
\end{quote}

The two nonemptinesses are fine --- though they come from the intersection
property applied to the BA-and-central $D\subseteq E$, not from the theorem
cited. \textbf{But ``hence'' does not follow}: two sets each meeting $E$ need
not meet each other. The maximality of $\sigma$, which would be the natural
extra ingredient, is not used anywhere in Case~1.

\paragraph{Why the source's Case~1 cannot be patched in place.}
The configuration \emph{reproduces itself} on $D$: one has $D\lll\bar A$,
$\bar C_{1}\cap D\lll D$, $\bar B_{2}\cap D\lll D$, still disjoint, with
$D\cap\bar B_{2}\neq\varnothing$. Every consequence of the two facts survives
the descent, so no iteration of them closes the case. The descent is not
useless --- $C_{1}\cap D\le_{\MT T}^{A}D$ properly, so the hypotheses are
inherited and an induction on $|B_{1}|$ is available --- but it shrinks the
ambient, while the conclusion's congruence family is indexed by the ambient
($\cover\omega\supseteq B_{1}^{2}$) and the single call site
(\texttt{main:2685}) consumes the original $B_{1}$ in one of its two branches.

\paragraph{The published repair, and why it fails.}
An earlier draft of this document passed to a \emph{minimal} subfamily
$F\subseteq\{\bar C_{1}^{1},\dots,\bar C_{1}^{t}\}$ with $\bigcap F\cap
\bar B_{2}=\varnothing$, ``which exists because the whole family works'', and
then used the stable-intersection theorem --- whose leave-one-out hypothesis is
exactly minimality, and whose four conclusions have no room for $\TS$ --- to
show no member of $F$ is $\TS$-typed, so that the case evaporates rather than
being contradicted.

The first line is wrong, and an external review found it. What the hypothesis
gives is $\bigl(\bigcap_{j}C_{1}^{j}\bigr)/\sigma\cap\bar B_{2}=\varnothing$,
whereas $\bigl(\bigcap_{j}C_{1}^{j}\bigr)/\sigma\subseteq
\bigcap_{j}\bigl(C_{1}^{j}/\sigma\bigr)$ with the inclusion generally strict.
So $\bigcap F$ can acquire $\sigma$-classes that meet $\bar B_{2}$, and the
subfamily need not exist. This is precisely the image-versus-intersection
failure that \S\ref{sec:leg-quotient} is about --- catalogued in this document
and then walked into. The lesson is recorded rather than tidied away.

\paragraph{Where it stands.}
\emph{Two things have changed since the repair was withdrawn.}

First, the image-versus-intersection failure is now known \emph{not} to occur
among the type-$T$ factors. If $\bar C_{1}^{j}\subt{T}\bar B_{1}$ then the
``moreover'' clause of \texttt{LEMFactorByDelta} gives
$\sigma\cap B_{1}^{2}\subseteq\sigma_{j}\cap B_{1}^{2}$, and a two-line
argument then shows
$\bigcap_{j\in F}\bar C_{1}^{j}=\bigl(\bigcap_{j\in F}C_{1}^{j}\bigr)/\sigma$
for $F$ the set of such $j$. So the trap that killed the published repair is
confined to the $\subt{\TS}$ and $=\bar B_{1}$ factors.

Second, given that, everything else in the argument runs: the proof paper
proves \pkey{lem:maximal-mult} in full from the single extra hypothesis
$(P\circ\sigma)\cap B_{2}=\varnothing$, where
$P=\bigcap_{j\in F}C_{1}^{j}$. No minimal subfamily, no $\TS$-freeness split,
and the ambient stays $\bar B_{1}$, which is what the conclusion needs.

So the whole of $\mathrm D17$ is now one sentence: \emph{the type-$T$ factors
alone already cut $B_{2}$ out}. That is strictly stronger than what the other
hypotheses give, since $C_{1}\subseteq P$.


The saturation identity $\delta\cap B^{2}\subseteq\sigma_{j}\cap B^{2}$ that
makes images commute with intersection is the ``moreover'' clause of the
factorisation lemma, and it is available for a factor precisely when that
factor lands in the type-$T$ alternative --- not when it lands in $\subt{\TS}$.
So the $\TS$ factors are the obstruction both to clause (m) of propagation and
to the minimal-subfamily argument, and the argument is circular as written.

One fragment survives, and it points at the target. A factor with
$\bar C_{1}^{j}\subt{\TS}\bar B_{1}$ contains a BA-and-central $\bar D$ of
$\bar B_{1}$, which meets $\bar B_{1}\cap\bar B_{2}$ by the strong-nonemptiness
lemma; so $\bar C_{1}^{j}\cap\bar B_{2}\neq\varnothing$ --- \emph{the $\TS$
factors never help cut $\bar B_{2}$ out}. What is missing is that the type-$T$
factors alone already cut it out. If that is proved, saturation applies to
them, the intersection of their images is the image of their intersection, and
the rest of the argument runs with the ambient still $\bar B_{1}$.

\begin{entry}[$\mathrm D17$ against the JACM proof]\label{ent:d17-jacm}
Checked on 2026-08-02 against the source of \cite{ZhukJACM}
(arXiv:1704.01914): \textbf{the lemma has no ancestor there}. The construction ---
a maximal congruence $\sigma$ with $(C_{1}\circ\sigma)\cap B_{2}=\varnothing$,
decomposed as an intersection of type-$T$ congruences --- appears nowhere in
the JACM paper, and the search is easy to reproduce: no occurrence of
``maximal congruence'' in its proof files does this job. The reason is
structural. There, a PC subuniverse \emph{is by definition} an intersection of
blocks of PC congruences, and a linear subuniverse \emph{is by definition}
stable under the minimal linear congruence --- the connection between
subuniverse-level and congruence-level data that
\texttt{LEMMaximalMultExtention} has to prove is definitional. The simplified
paper's chain-of-$\TD$-steps definitions buy the six-type uniformity that
\S5's calculus runs on, and this lemma is where the bill arrives. So the gap
cannot be closed by citation; it is an artifact of the new definitions, and it
is the simplified paper's own obligation.

The comparison also yields the first live repair route (the fourth recorded
in the proof paper's Caveat on the hypothesis): the JACM proof of the
parallelogram property proceeds by recursion over a measure on reductions,
restricting to the minimal linear reduction containing a solution of the
weakened instance (\texttt{ParPropertyMain}), with
\texttt{FitInLinearSubuniverses} --- chains of one-of-four subuniverses
through a subdirect relation intersect --- doing the work that stable
intersection does here. No maximal congruence appears. Since
\pkey{lem:maximal-mult} has exactly one consumer, re-proving
\pkey{lem:parallelogram-crucial} along those lines would delete the call site
and the hypothesis with it.

Two further facts from the same reading. The JACM induction measure is
$\operatorname{size}(D)=(|D_{1}|,\dots,|D_{s}|)$ over the \emph{distinct}
domains in decreasing order, compared lexicographically, with the explicit
remark that
duplicating domains does not affect it --- so the JACM theorems \emph{are}
applied at expanded coverings, and the measure is built to survive that. The
reviewer's \S3 objection to our $\mu$ describes the JACM architecture; the
simplified paper restructured the induction to appeal only at weakenings,
which is why $\mu$ suffices there (Entry~\ref{ent:measure}). And the JACM
route to the parallelogram property never types the $\ConOne$ congruences ---
the typing conclusion of \pkey{lem:parallelogram-crucial} is likewise new to
the simplified paper, so the fourth route must recover it separately:
irreducibility from \pkey{lem:crucial-irreducible}, the type from bridges.
\end{entry}

Proof paper: \pkey{lem:maximal-mult}, marked \emph{open} in its status table;
Caveat \pkey{haz:maximal-mult} now records the fourth route beside the three
failed ones.

\begin{entry}[Route four executed: what closed, and what the residue
  is]\label{ent:route-four}
Executed on 2026-08-02. The proof paper now proves
\pkey{lem:parallelogram-crucial} without citing
\pkey{lem:maximal-mult} at all, so Hypothesis \pkey{hyp:maximal-mult} is
consumed by nothing and $\mathrm D17$ no longer blocks the chain. The
residue is one claim and one stated lemma, and three findings from the
execution are worth recording.

\emph{First, the typing does not need the maximal-congruence apparatus in
the generic case.} With $E$ the trace of the crucial constraint at the
coordinate and $\tau_{1},\dots,\tau_{t}$ the witnessing congruences of the
multi-type reduction, cruciality applied to the weakened constraint $\exists
z\,R(z,x_{2},\dots,x_{n})\wedge\tau(z,x_{1})$, for $\tau$ a congruence
containing $\ConOne$, yields a clean dichotomy: either $\tau\subseteq\ConOne$
--- and then, because an irreducible congruence's cover is the \emph{least}
stable subuniverse properly above it, $\ConOne$ \emph{equals} a $\tau_{j}$
whenever a single $\tau_{j}\subseteq\ConOne$, which types it outright --- or
the weakening's solution produces a witness in $E$ inside the $\tau$-blocks.
The whole difficulty is confined to the second branch.

\emph{Second, the second branch cannot be closed at the level of subsets of
one algebra, which is a post-mortem for all three failed routes of
$\mathrm D17$.} In $\mathbf Z_{p}\times\mathbf Z_{p}$ the three coset
congruences of $\langle(1,0)\rangle$, $\langle(0,1)\rangle$,
$\langle(1,1)\rangle$ are linear and irreducible with full covers, and three
of their blocks meet pairwise with empty triple intersection; every
set-level datum of the bad branch is realized, no contradiction is
available, and the typing conclusion happens to be true there for reasons
the set-level data do not witness. Any proof must use that $E$ is the trace
of a crucial \emph{relation} --- exactly the hypothesis the failed routes
never used, and exactly why the JACM engine
(\texttt{FitInLinearSubuniverses}, \S\ref{sec:D17}'s Entry above) carries
tuple witnesses through the relation rather than sets alone.

\emph{Third, the source's own \TeX{} contains, commented out, a lemma of
exactly the shape the residue needs.} At \texttt{main:1851--1862} the v2
source carries (inside a \TeX{} comment, hence no part of the published
argument):

\begin{quote}\small
``Suppose $B\lll A$, $B$ is S-free, $T\in\{PC,L\}$, $\Sigma$ is the set of
all congruences on $A$ of type $T$ such that $\sigma^{*}\supseteq B^{2}$ and
$\sigma\not\supseteq B^{2}$, $\delta$ is an irreducible congruence on
$\mathbf A$ such that $\bigcap_{\sigma\in\Sigma}\sigma\subseteq\delta$ and
$B^{2}\not\subseteq\delta$. Then $\delta\in\Sigma$.''
(\texttt{LEMConOneIsPcOrLinear}, commented out)
\end{quote}

The proof paper states this, for a finite family, as
\pkey{lem:typed-above}, and its good branch consumes it; the $t=1$ case is
proved there. Note the commented lemma carries \emph{``$B$ is S-free''} ---
corroborating that the $\TS$-freeness hypothesis restored by $\mathrm D20$
is the right price, and consistent with the observation that
$\TS$-degradation cannot occur in quotients of $\TS$-free algebras (pull the
absorbing-and-central witness back along the quotient map, as in
\pkey{lem:reverse-hom}(bt)), which confines the old Hypothesis's failure
modes to collapsed factors.

Residue: \pkey{cl:typing-anchor} --- the bad branch, the type-$\TD$
Case-B corner, and $n=2$ for $T=\TPC$ --- and \pkey{lem:typed-above} for
$t\ge2$. Both are recorded as open debts in the proof paper's status table;
$\mathrm D19$ is the corresponding defect of the source's own argument at
the same site. Since this entry was written, the $\TBA$/$\TC$ Case-B
corner has come \emph{off} the residue:
\pkey{lem:anchored-intersection}, the anchored statement called for in
\S\ref{sec:D19}, excludes that configuration outright, and the claim now
carries only the remainder.
\end{entry}

\subsection{$\mathrm C{10}$ --- a citation that drops a case}\label{sec:C10}

\texttt{SS:208} states:

\begin{quote}\small
``Suppose $R\sdle\mathbf A\times\mathbf B$, $C=\{c\in A\mid\forall b\in B:
(c,b)\in R\}$. Then one of the following holds: (1) $C$ is a central
subuniverse of $\mathbf A$; (2) $\mathbf B$ has a nontrivial binary absorbing
subuniverse.''
\end{quote}

\cite[Thm.~6.15]{ZhukStrong} (p.~38) is verbatim this, except for a third item:
``3. $\mathbf B$ has a nontrivial \emph{projective} subuniverse.'' The word
``projective'' does not occur anywhere in the 2024 source ---
\texttt{grep -i projective} over all four of its \TeX{} files returns nothing
--- so the notion was removed from the paper, not merely from this citation.

\emph{The two-case statement is true, and the elimination is an argument Zhuk
himself runs elsewhere without ever connecting it to Theorem~6.15.} The
standing hypothesis (\texttt{main:1123}, restated at \texttt{main:1641}
precisely for the statements \S5 proves) makes $\mathbf B$ Taylor. By
\cite[Lem.~3.4]{ZhukStrong}, a nontrivial projective subuniverse that is not
binary absorbing produces an essentially unary $\mathbf U\in\mathrm{HS}(\mathbf
B)$ of size $\ge2$ --- which a WNU forbids. Zhuk runs this same move on this
same trichotomy in the proof of his Theorem~4.15, step (4)$\Rightarrow$(5):
``In case (5) Lemma 3.4 gives us a nontrivial binary absorbing subuniverse,
which is a strong subuniverse, or an essentially unary algebra $U\in HS(A)$,
which contradicts condition (4).''

The elimination should be stated as a lemma with a proof rather than smuggled
into a citation, and the proof paper does that, with a direct two-line argument
for ``a WNU forbids an essentially unary quotient'' in place of the source's
route through WNU-blockers.

\emph{But the two-case form is true only because $\varnothing$ counts as
central.} Under our reading (\S\ref{sec:leg-empty}) it is false:
$\mathbf A=\mathbf B=\mathbf Z_{p}$ with $R$ the graph of $x\mapsto x+1$ is
subdirect, has $C=\varnothing$, and $\mathbf Z_{p}$ has no nontrivial binary
absorbing subuniverse at all. So a third alternative $C=\varnothing$ must be
added. It is harmless: each of the nine call sites exhibits an element of $C$
first.

Proof paper: \pkey{lem:no-unary-quotient}, \pkey{lem:central-relation},
\pkey{haz:central-relation}.

\begin{entry}[Two obligations at each of the eight contrapositive call sites]
\label{ent:c10-sites}
The call sites are \texttt{SS:268, 551, 606, 843, 928, 1008, 2024, 2079}. Eight
run the lemma contrapositively against a \emph{BA and center free} algebra,
defined at \texttt{main:1363} as having ``no proper nonempty binary absorbing
subuniverse or proper nonempty central subuniverse''. Two obligations are
therefore live at each and discharged at none explicitly.

\emph{Properness of the centre.} Nonempty central is not enough; $C\neq A$ must
be shown. It is available at every site, and at three of them the witness is
printed two lines earlier --- the non-containments at \texttt{SS:533, 824, 903}
--- but never connected to the appeal. Sites \texttt{606, 1008, 2079} instead
\emph{use} the improper case, concluding that $R$ is full; there the split is
explicit.

\emph{Power to base.} Alternative (2) yields a binary absorbing subuniverse of
a \emph{power}; the sites assert one on the base. That is
\texttt{LEMBACenterSOnPowerImplies}, cited at \texttt{SS:2024} and omitted at
\texttt{551, 843, 928}.
\end{entry}

\subsection{$\mathrm F$ --- a full fibre from a surjective projection}\label{sec:F}

Subsubcases 1B3 (\texttt{SS:545}), 2A3 (\texttt{SS:837}) and 2B3
(\texttt{SS:921}) each contain, verbatim modulo names:

\begin{quote}\small
``Since $\proj_{1,2,\dots,|A|}(R')=(B/\delta)^{|A|}$, there exists
$d\in B_{m-1}/\sigma_{m}$ such that $(B/\delta)^{|A|}\times\{d\}\subseteq R'$.''
\end{quote}

Surjectivity of a projection does not produce a full fibre, and as stated this
is a non sequitur.

\emph{It is nevertheless true, for a reason that explains the arity.} Each such
$R'$ is cut out by a conjunction of unary and binary conditions on its first
$|A|$ entries together with a condition relating each entry to the last
coordinate --- at \texttt{SS:837}, $\forall i,j:(a_{i},a_{j})\in\sigma\cap
\omega$, $\forall i:a_{i}\in B_{m-1}$, and $(a_{i},b)\in\sigma_{m}$ --- and any
such condition survives substituting entries for one another. And
$|B/\delta|\le|A|$ because $B\subseteq A$. So the tuple enumerating $B/\delta$
with repetition already lies in $(B/\delta)^{|A|}$, its witness $d$ serves, and
every other tuple is a coordinate-substitution instance of it.

\textbf{The exponent $|A|$ is load-bearing}: it must be at least the number of
blocks, so that one tuple can enumerate them. A rendering that normalises the
arity away breaks all three proofs. Proof paper: \pkey{lem:full-fibre},
\pkey{haz:full-fibre}.

\subsection{$\mathrm D{18}$ --- the definition of a linear congruence is ill-typed}
\label{sec:D18}

\texttt{main:1373}, clause (3) of the definition of a linear congruence, reads
verbatim:

\begin{quote}\small
``there exist prime $p$ and $S\le(\sigma^{*})^{[4]}$ such that for any block
$B$ of $\sigma^{*}$ there exists $n\ge0$ with
$(B/\sigma;S\cap(B/\sigma)^{4})\cong(\mathbb Z_{p}^{n};x_{1}-x_{2}=x_{3}-x_{4})$.''
\end{quote}

$S$ is a relation on \emph{elements} of $A$ --- it is a subalgebra of
$(\sigma^{*})^{[4]}\subseteq A^{4}$ --- while $(B/\sigma)^{4}$ is a set of
$4$-tuples of \emph{$\sigma$-classes}. The two do not intersect: they are sets
of different things, and $S\cap(B/\sigma)^{4}$ is empty on inspection, or
meaningless, depending on how literally one reads it.

\emph{Repair.} Cut down inside $A$ and then take the image:
$(S\cap B^{4})/\sigma$. That is well formed for any $S$, and it is plainly what
is meant, since the isomorphism is with a relation on $\mathbb Z_{p}^{n}\cong
B/\sigma$.

The alternative repair --- restrict the image relation $S/\sigma$ to
$(B/\sigma)^{4}$ --- gives the same relation, because $S$ is stable under
$\sigma$ in each variable and its tuples meet no other block. But that
stability is a \emph{consequence} of clause (3) (it is what makes $S$ a bridge,
the source's own following remark), so it is not available while clause (3) is
being read, and a definition may not lean on its own consequences. Proof paper:
\pkey{def:linear-congruence}, \pkey{haz:linear-typing}, and the propagation
into \pkey{lem:linear-bridge}.

Note that the corresponding clause of \texttt{LEMNontrivialReflexiveBridgeImplies}
is typed \emph{correctly} in the source --- it writes
$(B/\sigma;(\delta\cap B^{4})/\sigma)$ --- so the slip is local to the
definition.

\subsection{$\mathrm D{19}$ --- two stable-intersection applications with no
  leave-one-out. \textbf{Open}}\label{sec:D19}

\texttt{CORMainStableIntersection} (\texttt{main:1803}) carries, as its
condition (4), the leave-one-out hypothesis: $R\cap(C_{1}\times\dots\times
B_{j}\times\dots\times C_{n})\neq\varnothing$ \emph{for every} $j$. The
typing argument of \texttt{LEMParalPropertyFromCrucialInMultiType} applies
the corollary twice, and neither application establishes condition (4), or
names the family it is applied to. The first, \texttt{main:2657--2666},
reads verbatim:

\begin{quote}\small
``If $E\cap D_{x_1}^{(1)}=\varnothing$, choose $C$, $B$, and $\mathcal
T\in\{\TBA,\TC,\TS,\TL,\TPC,\TD\}$ such that $D_{x_1}^{(1)}\lll^{D_{x_1}}
C<_{\mathcal T}^{D_{x_1}}B \lll D_{x_{1}}$, $E\cap C=\varnothing$, and
$E\cap B\neq \varnothing$. Since $R^{(1)}$ is not empty and $D^{(2)}\le_{
\mathcal M T}^{D}D^{(1)}$, Corollary [\texttt{CORMainStableIntersection}]
implies that $\mathcal T = T$.''
\end{quote}

$E\cap B\neq\varnothing$ is leave-one-out for the coordinate-$1$ member
only, and ``$R^{(1)}$ is not empty'' is nonemptiness at \emph{all} the
$B_{j}$ simultaneously, which is neither condition (3) nor condition (4).
The second application, \texttt{main:2693--2704}, concludes a critical
\emph{pair} --- ``there exist $i,j\in[n]$ [and PC blocks] such that $R$ has
no tuples whose $i$-th element is from $B_{i}$ and $j$-th element is from
$B_{j}$'' --- whose unstated derivation must shrink the family of all block
restrictions to a critical subfamily.

The standard repair --- shrink to a minimal inconsistent subfamily, so that
minimality supplies condition (4) --- stalls at a structural point: a
dropped member reverts its coordinate to the \emph{full} domain, but the
inconsistency being fed in ($E\cap C=\varnothing$, i.e.\ $R$ empty on a box
with $D^{(2)}$ at coordinates $2,\dots,n$) restricts those coordinates to
$D^{(2)}$ --- and at a coordinate where $D^{(2)}_{x}=D^{(1)}_{x}$ the
restriction $D^{(1)}_{x}$ is a $\lll$-chain, not a block of a typed
congruence, so it cannot be a member of the family at all. The
inconsistency the shrinking needs, with full domains at such coordinates,
is not known. What repairs it is an \emph{anchored} form of the
corollary --- known nonemptiness at the reduction rather than at the full
domains, which is exactly what $1$-consistency of $D^{(1)}$ supplies ---
a sibling of \texttt{FitInLinearSubuniverses} of the JACM proof rather
than of the corollary as printed.

\emph{Status: partly repaired.} The proof paper now carries that anchored
form: \pkey{lem:anchored-\brk intersection}, whose proof manufactures the
leave-one-out configurations by walking the multi-type chain to the first
step at which the intersection dies, and feeds that critical pair to the
relational corollary. It settles every application in which the two types
in the resulting pair agree with \emph{no} alternative of the corollary ---
in particular the first application above whenever $\mathcal
T\in\{\TBA,\TC,\TS\}$ or $\mathcal T$ is of the class opposite to $T$,
which in the proof paper's rewriting is the exclusion of configuration
$(\gamma)$ in the proof of \pkey{lem:parallelogram-crucial}
(Entry~\ref{ent:route-four}). What remains open is the \emph{same-class}
corner, where an alternative of the corollary is consistent with the data:
the first application with $\mathcal T=T$ needs no repair (it asserts
exactly $\mathcal T=T$), but the second application --- the critical
\emph{pair} of $\TPC$ blocks, and with it $n=2$ for $T=\TPC$ --- still
wants an anchored statement for a family of steps of one class, and that
is carried by \pkey{cl:typing-anchor}.

\subsection{$\mathrm D{20}$ --- the parallelogram lemma omits hypotheses its
  proof spends}\label{sec:D20}

\texttt{LEMParalPropertyFromCrucialInMultiType} (\texttt{main:2574})
hypothesises only: a $1$-consistent reduction $D^{(1)}$ for a constraint
$R(x_{1},\dots,x_{n})$, $T\in\{\TL,\TPC,\TD\}$, $D^{(2)}\le_{\mathcal
MT}^{D}D^{(1)}\lll D$, and $R$ crucial in $D^{(2)}$. Nothing makes $R$
subdirect over the \emph{unreduced} domains. Its own proof spends
subdirectness twice: at \texttt{main:2651}, ``Hence $E\neq\varnothing$ and
by Corollary [\texttt{CORPropagateToRelations}](r1) $E\lll D_{x_1}$'' ---
(r1) requires a subdirect $R$ with the full domains as ambient --- and in
the conclusion itself, where $\ConOne(R,i)$ can only be a congruence on the
full domain $\mathbf D_{x_{i}}$ if it is reflexive there, i.e.\ if
$\proj_{i}(R)=D_{x_{i}}$.

The same omission recurs at
\texttt{LEMMinimalPCLinearReductionIsConsistent} (\texttt{main:2262}): its
proof opens with ``By Lemma [\texttt{LEMPropagation}](fm)
$\proj_{i}(R^{(2)})\le_{\TMD}^{D_{x_{i}}}D_{x_{i}}^{(1)}$'', an application
whose ambient is the full domains, while the lemma hypothesises only that
the \emph{reduction} $D^{(1)}$ is $1$-consistent. (That proof also
concludes in $\le_{\TMD}$ and appeals to minimality among $\TMD$
subuniverses where its statement supplies minimality among
$\mathcal{M}T$ ones; propagation preserves the type, so the slip is
notational, and the proof paper's version keeps $T$ throughout.)

\emph{Repair.} Restore the hypothesis: the instance is
$1$-consistent. Both consumers supply it, their instances being
cycle-consistent. The proof paper's restatement also adds $\TS$-freeness of
the $\mathbf D^{(1)}_{x_{i}}$, which is not spent by the source's own route
but is what the replacement route's \pkey{lem:typed-above} carries --- and
which the source's commented-out \texttt{LEMConOneIsPcOrLinear} also
carries (Entry~\ref{ent:route-four}). At the consumer inside the main
induction the freeness is automatic (no domain there has a nontrivial
binary absorbing or central subuniverse); at the codimension-one consumer
it is a new, recorded obligation on the blocks $\phi(\bar b)_{x}$. Proof
paper: \pkey{lem:parallelogram-crucial}, hypotheses (i) and (ii), and its
Caveat \pkey{haz:parallelogram-hyp}; \pkey{lem:minimal-consistent}, now
proved with the restored hypothesis, and its Caveat
\pkey{haz:minimal-consistent}.

\subsection{$\mathrm D{21}$ --- the linkedness-preservation auxiliary asserts
  its cross-component step. \textbf{Open}}\label{sec:D21}

\texttt{LEMPreserve\brk Linkdness\brk OneStepAUX} (\texttt{SS:2821--2896})
is the workhorse behind \texttt{LEMPreserve\brk Linkdness}; the chain is AUX
$\to$ \texttt{LEMPreserve\brk Linkdness\brk OneStep} (\texttt{SS:2984})
$\to$ the lemma
(\texttt{SS:3007}), and the outer two links are sound
(\S\ref{sec:exp-preserve}). Case~2 of AUX's own proof ends, verbatim
(\texttt{SS:2892--2895}):

\begin{quote}\small
``Since $S\cap(C_{1}\times C_{2})\neq\varnothing$,
$C_{1}'\times C_{1}''\subseteq \mathrm{LeftLinked}(R')$, which contradicts
condition 1 of Lemma [\texttt{LEMCongruenceEitherCutOrDoNothing}].''
\end{quote}

No justification is given, and none is visible. $S\cap(C_{1}\times
C_{2})\neq\varnothing$ places one pair in one linked component of $R$;
$C_{1}''$ is produced by a \emph{different} nonemptiness,
$R'\cap(B_{1}'\times C_{2}')\neq\varnothing$, and nothing places its
witness in the same component --- $R'$ could a priori split into components
with the $S$-witness in one and the mixed edge in another. The saturation
$\mathrm{LeftLinked}(R')\supseteq\delta_{1}$ is true and useful but joins
only $\delta_{1}$-related points, and $C_{1}'$ is a union of several
$\delta_{1}$-blocks.

The gap is exactly a \emph{component-attachment} problem, and it is the
whole difficulty of the lemma: in the regime where Case~2 is entered,
condition 1 of \texttt{LEMCongruence\brk Either\brk Cut\brk Or\brk
DoNothing} pins every linked
component to a single block, so the components carry disjoint private
matchings, and the printed hypotheses do not connect the component of the
$S$-witness to any edge in the box. The fullness regime, by contrast, is
closed by \texttt{LEMLeftLinkedStayFull} unconditionally --- so the defect
is confined to, and coextensive with, the cut regime.

Three further facts, established on the proof-paper side. First, the
statement survives one strengthening: with an \emph{anchor} --- every
element of $C_{1}$ has an $R$-neighbour in $B_{2}$ --- the witness's own
neighbour is an edge of the right component, and the full/cut dichotomy
closes both ways; this is \pkey{lem:preserve-linked-anchored},
\textsc{proved}. Second, the anchor is not dischargeable at the lemma's one
consumer (\texttt{main:2632}, the parallelogram step of
\texttt{LEMParalPropertyFromCrucialInMultiType}): there it would say that
every section of the crucial constraint over one block of coordinates with
entries in $D^{(2)}$ extends through $R$ to the other block inside
$D^{(1)}$, and $1$-consistency of $D^{(1)}$ is per-variable, not
per-tuple. Third, no counterexample can live in the purely affine world:
a subdirect subuniverse of a product of affine algebras is a coset, and a
coset is its own rectangular closure, so there the conclusion holds
outright. The source itself contains, commented out at
\texttt{SS:2899--2981}, an abandoned alternative proof of a one-sided
variant carrying the extra hypothesis
$R\cap(B_{1}\times B_{2})\le_{sd}B_{1}\times B_{2}$ --- an anchor ---
which suggests the difficulty was met and sidestepped.

\emph{Status: open.} \pkey{lem:preserve-linked} remains stated;
\pkey{lem:preserve-linked-anchored} is its proved core, and Caveat
\pkey{haz:preserve-anchor} records the obstruction. The live repair route
is not to prove the unanchored lemma but to bypass it at the consumer:
cruciality reduces the parallelogram step to a single zigzag
($W\circ W^{-1}\circ W\supsetneq W$ would itself be a properly weaker
constraint), and the per-coordinate descent that remains is of the same
anchored shape as \pkey{lem:find-consistent} and
Claim~\pkey{cl:typing-anchor} --- folding this defect into the debts the
chain already carries rather than adding a new one.

\section{Statements restated}\label{sec:restated}

Five statements of the source are restated in the proof paper, not because they
are wrong but because our conventions make the printed form say something else
--- or because the source's own proof gives more than its statement.

\begin{center}\footnotesize
\begin{tabular}{@{}p{0.24\textwidth}p{0.32\textwidth}p{0.34\textwidth}@{}}
\hline
\textbf{Statement} & \textbf{Why} & \textbf{Restatement} \\
\hline
\texttt{LEMCentral\brk RelationImplies} (\texttt{SS:208}) &
  True in the source only because $\varnothing$ counts as central, and the
  projective alternative is dropped (\S\ref{sec:C10}) &
  ``$C=\varnothing$, or $C$ is central in $\mathbf A$, or $\mathbf B$ has a
  proper nonempty binary absorbing subuniverse'' \\
\texttt{LEMLinked\brk ImpliesBACenter} (\texttt{SS:222}) &
  ``There exists a BA or central subuniverse'' is vacuous: every algebra is a
  binary absorbing subuniverse of itself &
  Insert \emph{proper nonempty}; that is how it is used at \texttt{SS:1282} \\
The $\subt{\TS}$ clause (\texttt{main:1522}) &
  Vacuously true for $D=\varnothing$ &
  Require the witness $D$ nonempty \\
\texttt{LEMIntersectALL}(i) &
  Its proof (\texttt{SS:2356}) establishes $B\cap D\dlll^{A}D$ and then weakens
  it; three later proofs use the stronger form while citing (i) &
  State (i) as $B\cap D\dlll^{A}D$ and derive $B\cap D\dlll A$ as a corollary \\
\texttt{CORReverseHomomorphism} &
  Its type list is $\{\TBA,\TC,\TS,\TL,\TD\}$ where the stable-intersection
  theorem uses $\{\TBA,\TC,\TS,\TL,\TPC\}$ &
  Fix one convention; harmless, since $\TD$ subsumes $\TL$ and $\TPC$ \\
\hline
\end{tabular}
\end{center}
